\chapter{万有引力与重力}\label{ch:g10:universal-gravitation}

松手放开一只苹果，它便下落。月球距地心六十个地球半径之遥，却似乎从不下落。
牛顿的伟大洞见在于，月球\emph{确实}在下落——
每秒约 1.4 毫米，恰等于
引力随距离平方减弱所预言的量。一个公式，配以一个
普适常量，支配着苹果、月球、潮汐与行星。
第一分册以文字描述了引力；
本章则要测量它。

\section{万有引力定律}

\begin{definition}[万有引力]\label{def:g10:universal-gravitation:force}
任意两个物体都相互吸引，无论它们由何种物质构成，
也无论它们相距多远。这种吸引称为\emph{万有引力
}\index{万有引力}，而这一现象称为
\emph{万有引力}\index{万有引力}。它无需接触、绳索，
也无需介质，并且它是\emph{相互的}：每个物体都拉动对方。
\end{definition}

\begin{theorem}[万有引力定律]\label{thm:g10:universal-gravitation:law}
质量分别为 $m_1$ 和 $m_2$（单位为 \unit{kg}）的两个物体，其中心相距
距离为 $d$（单位为 \unit{m}），两者以两个力相互吸引，
两力沿两中心连线方向，各自指向另一物体，
大小相等，
\[
F = G\,\frac{m_1 m_2}{d^2}
\qquad \text{（单位为 \unit{N}）},
\]
其中 $G$ 对宇宙中每一对物体都是同一常量。
\end{theorem}
\admitted

\begin{remark}[定律的来源]
牛顿由观测到的行星运动归纳出这一定律
（即开普勒定律）；严格的推导在
本科一年级分册中给出。\Cref{ch:g12:satellites-kepler}已将
该定律用于卫星与行星。
\end{remark}

\begin{definition}[引力常量]\label{def:g10:universal-gravitation:constant}
该常量
\[
G = \qty{6.67e-11}{N.m^2/kg^2}
\]
即\emph{引力常量}\index{引力常量}。其
数值极小，正是日常物体间的引力
难以察觉的原因：只有至少一个质量达到天文量级，
$G m_1 m_2 / d^2$ 才有意义。
\end{definition}

\begin{omfigure}
\begin{tikzpicture}[scale=1.0]
  \draw[omExo, dashed] (0,0) -- (5.2,0);
  \fill[omDef!25] (0,0) circle (0.45);
  \draw[omDef, thick] (0,0) circle (0.45);
  \fill[omDef!25] (5.2,0) circle (0.28);
  \draw[omDef, thick] (5.2,0) circle (0.28);
  \node[font=\small] at (0,0) {$m_1$};
  \node[font=\small] at (5.2,0) {$m_2$};
  \draw[-Stealth, omThm, very thick] (0.45,0) -- (1.55,0)
    node[midway, above, font=\small] {$\vect F_{2/1}$};
  \draw[-Stealth, omThm, very thick] (4.92,0) -- (3.82,0)
    node[midway, above, font=\small] {$\vect F_{1/2}$};
  \draw[Stealth-Stealth, thin] (0,-0.75) -- (5.2,-0.75)
    node[midway, below, font=\small] {$d$};
\end{tikzpicture}

{\small 引力对中的两个力：沿中心连线方向，
指向另一物体，大小\emph{相同}，
$F = G m_1 m_2 / d^2$ 即使两物体质量相差悬殊，亦然。}
\end{omfigure}

\begin{remark}[读懂公式]\label{rem:g10:universal-gravitation:reading}
公式有四个值得注意的特点。
\begin{itemize}
  \item \emph{相互且等大。}地球以你的重力拉你，
        ——你也以完全相同的力拉地球。（地球
        只是反应较小：它的质量约为你的
        $\num{e23}$ 倍。）
  \item \emph{平方反比。}$d$ 加倍，力便除以 $4$；
        距离扩大到十倍，力便除以 $100$。
  \item \emph{中心距离。}对于球形物体——恒星、
        行星、炮弹——$d$ 应从中心量到中心，而非
        从表面量到表面。
  \item \emph{两质量皆重要。}任一质量加倍，力即加倍；
        公式对 $m_1$ 与 $m_2$ 对称。
\end{itemize}
\end{remark}

\begin{example}[两个朋友]\label{ex:g10:universal-gravitation:friends}
两位质量为 \qty{70}{kg} 的朋友，中心相距 \qty{1.0}{m}：
\[
F = \num{6.67e-11} \times \frac{70 \times 70}{1.0^2}
\approx \qty{3.3e-7}{N}.
\]
每位朋友的重力约为 \qty{6.9e2}{N}——
是两人间引力的\emph{二十亿}倍。两人间的引力确实存在，却无从感知。
\end{example}

\begin{example}[地球束缚月球]\label{ex:g10:universal-gravitation:earthmoon}
取 $m_1 = \qty{5.97e24}{kg}$（地球），$m_2 = \qty{7.35e22}{kg}$（月球）
及 $d = \qty{3.84e8}{m}$：
\[
F = \num{6.67e-11} \times
\frac{\num{5.97e24} \times \num{7.35e22}}{(\num{3.84e8})^2}
\approx \qty{2.0e20}{N}.
\]
月球以相同的 \qty{2.0e20}{N} 拉力拉回地球——
这一拉力使海洋每天两次作出响应，形成潮汐。
\end{example}

\begin{method}[计算万有引力]\label{met:g10:universal-gravitation:compute}
\begin{enumerate}
  \item 将质量换算为千克，将两中心
        之间的距离换算为米。
  \item 应用 $F = G m_1 m_2 / d^2$。
  \item 将 10 的幂与数字部分分开检查，
        如\cref{ch:g10:orders-of-magnitude}所示，万有引力的取值范围
        从 $\num{e-7}$ 牛顿（两人相距一米）到 $\num{e22}$ 牛顿
        （太阳对地球），因此指数写错很容易发现。
\end{enumerate}
\end{method}

\section{重力与重力加速度}

\begin{definition}[重力]\label{def:g10:universal-gravitation:weight}
\emph{重力}\index{重力} $\vect P$ 是物体在恒星
（行星、卫星等）表面所受恒星施加的万有引力。
其方向指向恒星中心，也就是我们所称的
\emph{竖直方向}\index{竖直方向}的方向，单位为牛顿。
\end{definition}

\begin{proposition}[恒星表面的重力]\label{prop:g10:universal-gravitation:surfaceg}
质量为 $M$、半径为 $R$ 的球形恒星表面，
质量为 $m$ 的物体所受重力为
\[
P = m\,g,
\qquad\text{其中}\qquad
g = \frac{G M}{R^2}
\]
仅取决于恒星，与物体无关。
\end{proposition}

\begin{proof}
物体中心与恒星中心相距 $R$，因此
由万有引力定律
（\cref{thm:g10:universal-gravitation:law}）可得
$P = G M m / R^2 = m \times \left(G M / R^2\right)$。括号内的量
对表面上的每个物体都相同，记为 $g$。
\end{proof}

\begin{definition}[重力加速度]\label{def:g10:universal-gravitation:fieldstrength}
量 $g = G M / R^2$，单位为牛顿每千克，是
恒星表面的\emph{重力加速度}\index{重力加速度}：
即恒星对每千克物质的引力为 $g$ 牛顿。
在地球上，$g = \qty{9.81}{N/kg}$。
\end{definition}

\begin{example}[计算 $g$]\label{ex:g10:universal-gravitation:gvalues}
对地球（$M = \qty{5.97e24}{kg}$，$R = \qty{6.37e6}{m}$）：
\[
g = \frac{\num{6.67e-11} \times \num{5.97e24}}{(\num{6.37e6})^2}
= \frac{\num{3.98e14}}{\num{4.06e13}}
\approx \qty{9.81}{N/kg}.
\]
对月球（$M = \qty{7.35e22}{kg}$，$R = \qty{1.74e6}{m}$），同样计算
得 $g_{\text{月球}} \approx \qty{1.62}{N/kg}$，约为地球的六分之一：
月球的质量小得多，半径较小也无法弥补这一差距。
对火星，$g_{\text{火星}} \approx \qty{3.73}{N/kg}$。
\end{example}

\begin{omfigure}
\begin{tikzpicture}
\begin{axis}[omaxis, axis x line=bottom, axis y line=left,
  width=8.8cm, height=5.4cm, ybar, bar width=15pt,
  symbolic x coords={Moon, Mars, Earth, Jupiter}, xtick=data, xticklabels={月球, 火星, 地球, 木星},
  ylabel={$g$ (\unit{N/kg})}, ymin=0, ymax=29,
  nodes near coords,
  every node near coord/.append style={font=\small},
  enlarge x limits=0.16]
\addplot[draw=omDef, fill=omDef!25] coordinates
  {(Moon,1.62) (Mars,3.73) (Earth,9.81) (Jupiter,24.8)};
\end{axis}
\end{tikzpicture}

{\small 四个天体表面的重力加速度：同样的
\qty{1}{kg} 糖袋在月球上的重力为 \qty{1.6}{N}，
在木星上则为 \qty{25}{N}。}
\end{omfigure}

\begin{omfigure}
\begin{tikzpicture}[scale=1.0]
  \fill[omExo!15] (0,0) circle (1.5);
  \draw[omExo, thick] (0,0) circle (1.5);
  \fill (0,0) circle (1.5pt) node[below, font=\small] {$O$};
  \fill (90:1.5) circle (1.8pt) node[above, font=\small] {$A$};
  \draw[-Stealth, omThm, thick] (90:1.5) -- (90:0.8)
    node[pos=0.6, right, font=\small] {$\vect P_A$};
  \fill (200:1.5) circle (1.8pt) node[left, font=\small] {$B$};
  \draw[-Stealth, omThm, thick] (200:1.5) -- (200:0.8)
    node[pos=0.6, below, font=\small] {$\vect P_B$};
  \fill (330:1.5) circle (1.8pt) node[right, font=\small] {$C$};
  \draw[-Stealth, omThm, thick] (330:1.5) -- (330:0.8)
    node[pos=0.6, above left, font=\small] {$\vect P_C$};
\end{tikzpicture}

{\small “向下”是一个局部概念：在地球的每一点，
重力矢量都指向球心 $O$。对跖点上的两位行人
都直立站立，脚朝向彼此。}
\end{omfigure}

\begin{remark}[宇航员为何会蹦跳]\label{rem:g10:universal-gravitation:bounce}
一名质量为 \qty{120}{kg}（连同宇航服）的宇航员在地球上的重力为
$120 \times 9.81 \approx \qty{1180}{N}$，而在月球上仅为
$120 \times 1.62 \approx \qty{194}{N}$。习惯于
\qty{1180}{N} 负荷的肌肉突然只需承担六分之一的负载，因此
形成了著名的缓慢蹦跳式月面步态。质量——以及随之而来的
\emph{停止}或\emph{转向}所需的努力——都没有改变，这正是
月球宇航员频频摔倒的原因。
\end{remark}

\section{质量还是重力？天平与弹簧测力计}

日常语言说一袋面粉“重三千克”。物理学
则把这句话一分为二。

\begin{remark}[质量不是重力]\label{rem:g10:universal-gravitation:massweight}
\begin{itemize}
  \item \emph{质量} $m$ 度量物体所含物质的多少，
        单位为千克。它只取决于物体本身：在地球上、
        在月球上，还是漂浮于深空中，都一样。
  \item \emph{重力} $P = mg$ 是一种力，单位为牛顿，由
        \emph{恒星}施加于物体上，其大小随
        $g$ 而变 --- 远离任何恒星时则趋于零。
\end{itemize}
混淆二者在市场买卖中无伤大雅，在物理
问题中却会导致大错：千克不是力的单位。
\end{remark}

\begin{method}[天平与弹簧测力计]\label{met:g10:universal-gravitation:instruments}
\begin{itemize}
  \item \emph{天平}将物体与标准质量相比较：
        两盘所受的重力都按同一 $g$ 成比例变化，因此这种比较
        测量的是\emph{质量} --- 在任何恒星上结论都相同。
  \item \emph{弹簧测力计}测量拉伸弹簧的力：
        它示数为\emph{重力}，单位为牛顿，其示数随
        恒星而变。要求出质量，需将其除以当地的
        重力加速度：$m = P/g$。
\end{itemize}
\end{method}

\begin{example}[月球上的面粉]\label{ex:g10:universal-gravitation:flour}
一袋 \qty{3.0}{kg} 的面粉与 \qty{3.0}{kg} 的标准
质量块在地球\emph{和}月球上都能平衡。然而，弹簧测力计
在地球上示数为 $3.0 \times 9.81 \approx \qty{29}{N}$，
在月球上为 $3.0 \times 1.62 \approx \qty{4.9}{N}$。一位商人在
月球上“按牛顿”售卖面粉，就会多给出六倍的面粉 ---
交易与科学一样，依靠的是质量。
\end{example}

\section{绕地球下落：牛顿大炮}

将一块石头水平抛出：重力使其轨迹弯成弧线，它
落在几米外。抛得更快，落点就更远。牛顿的
思想实验把这一想法推到极限：地球是圆的，而
沿地表水平方向每前进 \qty{8}{km}，地表便低于水平线约 \qty{5}{m}；
而自由落体在第一秒内下落 \qty{4.9}{m}
（这一点将在\cref{ch:g12:projectile-motion}中证明）。因此，
以 \qty{8}{km/s} 速度贴地飞行的抛体，其下落
速度恰好等于地面弯离的速度：它将永远下落，却永远不会落地。

\begin{definition}[卫星与轨道]\label{def:g10:universal-gravitation:satellite}
称绕恒星下落、从不触及其表面的物体为该恒星的\emph{卫星}\index{卫星}；
它反复运行的路径称为它的
\emph{轨道}\index{轨道}。
\end{definition}

\begin{omfigure}
\begin{tikzpicture}[scale=1.0]
  \fill[omExo!15] (0,0) circle (1.5);
  \draw[omExo, thick] (0,0) circle (1.5);
  \fill[omExo!60] (-0.22,1.44) -- (0,1.85) -- (0.22,1.44) -- cycle;
  \coordinate (L) at (0,1.85);
  \draw[omThm, thick] (L) .. controls (0.6,1.85) .. (55:1.5);
  \draw[omProp, thick] (L) .. controls (1.1,1.85) and (1.7,1.1) .. (10:1.5);
  \draw[omMeth, thick] (L) .. controls (1.8,1.8) and (2.3,-0.4) .. (-70:1.5);
  \draw[-Stealth, omDef, thick] (0,1.85) arc (90:-250:1.85);
  \node[font=\small, omDef, anchor=east] at (-1.5,1.45) {轨道};
\end{tikzpicture}

{\small 牛顿大炮：从山顶发射，速度越来越快，炮弹
落点越来越远 --- 直到约 \qty{7.9}{km/s} 时，地面
弯离的速度恰与炮弹下落速度相同，炮弹便绕地球运行。}
\end{omfigure}

\begin{remark}[月球是一块下落的石头]\label{rem:g10:universal-gravitation:fallingmoon}
月球是地球的天然卫星：它每秒向地球下落约
\qty{1.4}{mm}，同时其侧向运动使它绕地球运行，
从而无法再靠近地球。
（在\cref{exo:g10:universal-gravitation:15}中可重做牛顿
本人对这一数值的检验）。绕轨道运行的空间站中，宇航员与飞船
处于相同的自由落体之中 --- 正是这一点，而非重力消失，导致他们漂浮。
引力塑造每一条轨迹，从
石头的弧线到行星的椭圆；轨道的定量理论将在
\cref{ch:g12:satellites-kepler}中讲述。
\end{remark}

\begin{notation}[数据卡]
除非练习另有说明，请取
$G = \qty{6.67e-11}{N.m^2/kg^2}$ 及以下数据：
\begin{center}
\small
\begin{tabular}{lccc}
 & 质量（\unit{kg}） & 半径（\unit{m}） & $g$（\unit{N/kg}） \\
地球 & \num{5.97e24} & \num{6.37e6} & 9.81 \\
月球 & \num{7.35e22} & \num{1.74e6} & 1.62 \\
火星 & \num{6.42e23} & \num{3.39e6} & 3.73 \\
太阳 & \num{1.99e30} & --- & --- \\
\end{tabular}
\end{center}
地球--月球距离：\qty{3.84e8}{m}；地球--太阳距离：
\qty{1.496e11}{m}（中心到中心）。
\end{notation}

\section{练习}

\begin{exercise}[$\star$]\label{exo:g10:universal-gravitation:1}
两名 \qty{60}{kg} 的舞伴相互拥抱，两人的中心
相距 \qty{0.80}{m}。
\begin{enumerate}
  \item 求两人之间的万有引力。
  \item 将其与一只 \qty{2.5}{mg} 蚊子的重力相比较。
\end{enumerate}
\end{exercise}

\begin{exercise}[$\star$]\label{exo:g10:universal-gravitation:2}
一名学生的质量为 \qty{55}{kg}。
\begin{enumerate}
  \item 求她在地球上的重力。
  \item 她前往月球。她在那里的质量和
        重力各是多少？
\end{enumerate}
\end{exercise}

\begin{exercise}[$\star$]\label{exo:g10:universal-gravitation:3}
\begin{enumerate}
  \item 根据数据卡，重新计算火星表面的
        重力加速度。
  \item 一辆 \qty{900}{kg} 的火星车被送往那里，求其在火星上的
        重力与在地球上的重力。
\end{enumerate}
\end{exercise}

\begin{exercise}[$\star$]\label{exo:g10:universal-gravitation:4}
根据数据卡，求地球对月球施加的
力。月球对地球施加的力是多少？
\end{exercise}

\begin{exercise}[$\star$]\label{exo:g10:universal-gravitation:5}
两物体之间以力 $F_0$ 相互吸引，则在下列情况下
该力变为多少：
\begin{enumerate}
  \item 两中心之间的距离加倍；
  \item 距离减半；
  \item 两质量均加倍（距离不变）；
  \item 一个质量变为原来的三倍，\emph{且}距离也变为原来的三倍？
\end{enumerate}
\end{exercise}

\begin{exercise}[$\star\star$]\label{exo:g10:universal-gravitation:6}
两艘满载的超级油轮各为 \qty{5.0e8}{kg}，系泊时
两船中心相距 \qty{100}{m}。
\begin{enumerate}
  \item 求两船之间的万有引力。
  \item 在地球上，重力等于此力的物体质量为多少？
  \item 该力大得惊人 --- 然而两船并无相互靠拢的
        迹象。求其给予一艘油轮的加速度
        $a = F/m$，并加以评论。
\end{enumerate}
\end{exercise}

\begin{exercise}[$\star\star$]\label{exo:g10:universal-gravitation:7}
一家月球商店售卖一袋 \qty{3.0}{kg} 的大米。
\begin{enumerate}
  \item 该袋大米在地球上和月球上，天平与弹簧测力计
        的示数分别是多少？
  \item 顾客按“\qty{3.0}{kg}”付款。哪种仪器
        能在任何星球上都保证交易公平？请说明理由。
\end{enumerate}
\end{exercise}

\begin{exercise}[$\star\star$]\label{exo:g10:universal-gravitation:8}
国际空间站在 \qty{400}{km} 的高度上飞行。
\begin{enumerate}
  \item 求地球在该高度处的重力加速度。（注意：
        此处的距离应为到地球
        \emph{中心}的距离吗？）
  \item 将其表示为地表处数值的百分比。
  \item 站上的宇航员处于漂浮状态。请将此现象与你的结果相协调。
\end{enumerate}
\end{exercise}

\begin{exercise}[$\star\star$]\label{exo:g10:universal-gravitation:9}
金星表面处 $g = \qty{8.87}{N/kg}$，该行星的半径为
\qty{6.05e6}{m}。求金星的质量，
并与地球的质量相比较。
\end{exercise}

\begin{exercise}[$\star\star$]\label{exo:g10:universal-gravitation:10}
\begin{enumerate}
  \item 求太阳对 \qty{1.0}{kg} 海水施加的力，
        然后求月球对同一千克海水施加的
        力。
  \item 太阳的拉力更大 --- 大了多少倍？
  \item 然而，月球主导潮汐。潮汐的驱动力并非拉力本身，
        而是地球近月一侧与远月一侧所受拉力之\emph{差值}。
        请不经计算说明，为何距离较近的月球
        在这一判据下能胜过距离较远的太阳，
        理由何在？
\end{enumerate}
\end{exercise}

\begin{exercise}[$\star\star$]\label{exo:g10:universal-gravitation:11}
\begin{enumerate}
  \item 某行星的质量是地球的两倍，半径也是地球的两倍。
        试用地球的 $g$ 表示其表面的重力加速度，
        再以 \unit{N/kg} 表示。
  \item 质量与地球相同的行星，若要使其表面
        重力加速度为 $2g$，半径应为多少？
\end{enumerate}
\end{exercise}

\begin{exercise}[$\star\star\star$]\label{exo:g10:universal-gravitation:12}
地球--月球连线上的某处，空间探测器所受的两个拉力
相互抵消。设 $x$ 为该点到地球中心的距离，
$D$ 为地球--月球距离。
\begin{enumerate}
  \item 证明在该点处，
        $\left(\dfrac{x}{D - x}\right)^{2}
        = \dfrac{M_{\text{地球}}}{M_{\text{月球}}}$。
  \item 由此求出 $x$。该距离占前往月球全程的几分之几？
\end{enumerate}
\end{exercise}

\begin{exercise}[$\star\star\star$]\label{exo:g10:universal-gravitation:13}
牛顿大炮，附数值计算。一抛体沿地球
水平方向掠过，速率为 $v$。
\begin{enumerate}
  \item 水平飞行 $x = \qty{8.0}{km}$ 后，球形地球
        因弯曲而“下降”了高度 $h$，满足
        $(R + h)^2 = R^2 + x^2$。利用
        $(R + h)^2 \approx R^2 + 2Rh$（$h$ 很小），证明
        $h \approx x^2 / (2R)$，并计算 $h$。
  \item 自由落体在第一秒内下落 \qty{4.9}{m}。
        将其与 $h$ 比较，并推出抛体
        永不落地的速率。
  \item 求以该速率贴地表运行时，绕行一周
        的时长。
\end{enumerate}
\end{exercise}

\begin{exercise}[$\star\star\star$]\label{exo:g10:universal-gravitation:14}
一颗中子星将质量 $M = \qty{2.8e30}{kg}$（为太阳质量的 1.4 倍）
压缩在半径为 \qty{12}{km} 的范围内。
\begin{enumerate}
  \item 求其表面的重力加速度。
  \item 一个质量为 \qty{70}{kg} 的人在那里的重力为多大？
        并与其在地球上的重力相比较。
  \item 一粒 \qty{1.0}{mg} 的沙粒在那里的重力是多少？
        在地球上，重力与之相同的物体质量是多少？
\end{enumerate}
\end{exercise}

\begin{exercise}[$\star\star\star$]\label{exo:g10:universal-gravitation:15}
牛顿的月球检验（1666 年）。月球轨道半径约为
$60$ 个地球半径。
\begin{enumerate}
  \item 若引力遵循平方反比律，则月球处
        每千克物体所受地球的引力，
        比地表处小多少倍？
  \item 在地表附近，一个自由下落的物体在
        第一秒内下落 \qty{4.9}{m}。月球每秒应
        向地球下落多远？
  \item 用轨道检验：月球绕半径为
        $d = \qty{3.84e8}{m}$ 的圆轨道一周需 $T = 27.3$ 天。求其速率
        $v$，其一秒内走过的距离 $x$，然后算出下落高度
        $h \approx x^2/(2d)$（方法同
        \cref{exo:g10:universal-gravitation:13}）。
  \item 比较第 2 题与第 3 题的结果，并说明牛顿
        得出了什么结论。
\end{enumerate}
\end{exercise}

\section{问题：称量地球}

\begin{problem}[{周末问题 --- 从棚屋中的两个铅球
到地球的质量，再到太阳的质量：一个只需测量一次的小常量，
究竟有何价值}]
\label{pb:g10:universal-gravitation:1}
1798 年，亨利·卡文迪什用细线悬挂起一根
携带两个小铅球的轻杆，将两个大铅球移近，
并测量了细线几乎难以察觉的扭转。报纸称他
已“\emph{称出了地球}”，报道并无夸大：一旦 $G$ 已知，
地球自身的 $g$ 与半径即可给出其质量，
地球一年的轨道则给出太阳的质量。本题将重走
这一整条推导链。全程使用数据卡。

\medskip
\noindent\textbf{第一部分 --- 最微弱的力。}
\begin{enumerate}
  \item 埃菲尔铁塔的铁质量约为 \qty{7.3e6}{kg}。一个
        \qty{55}{kg} 的游客站在距其质心 \qty{100}{m}
        处。求铁塔对该游客的引力。
  \item 将其与一粒 \qty{1.0}{mg} 沙粒的重力相比较。
  \item 比例练习：引力会如何变化？分别考虑
        (a) 距离变为原来的 $10$ 倍；(b) 两质量
        均变为原来的 $1000$ 倍；(c) 两质量\emph{和}距离
        均变为原来的 $1000$ 倍？
  \item 计算地球对位于
        地表处的 \qty{1.00}{kg} 物体的引力（由 $M$、$R$ 与 $G$ 求得）。这一数值有什么
        日常名称，用何种符号表示？
  \item 引力是你所学过的相互作用中最为微弱的 ---
        却支配着整个宇宙。用一两句话说明，
        为何地球对你的引力远大于对你邻居的引力。
\end{enumerate}

\medskip
\noindent\textbf{第二部分 --- 卡文迪什扭秤。}
在扭秤中，质量为 $m_1 = \qty{158}{kg}$ 的大球
吸引质量为 $m_2 = \qty{0.73}{kg}$ 的小球；两者球心
相距 $d = \qty{0.23}{m}$。
\begin{enumerate}[resume]
  \item 取 $G$ 的现代值，计算两球间的
        万有引力。
  \item 计算小球的重力，以及两个力的比值。卡文迪什为何
        需要扭丝，而不用
        任何普通的秤？
  \item 卡文迪什测得的扭转（以现代单位计）
        对应于 $F = \qty{1.47e-7}{N}$。推算他的 $G$ 值，
        并与今天的值比较。
  \item 取 $G = \qty{6.67e-11}{N.m^2/kg^2}$，$g = \qty{9.81}{N/kg}$，
        $R = \qty{6.37e6}{m}$，推算地球的质量（这
        正是报纸津津乐道的一步）。
  \item 推算地球的平均密度 $\rho = M / V$，其中
        $V = \frac43 \pi R^3$。地表岩石的密度约为
        \qty{2.7e3}{kg/m^3}，比较二者可说明地球深部的
        什么情况？
\end{enumerate}

\medskip
\noindent\textbf{第三部分 --- 其他星球。}
\begin{enumerate}[resume]
  \item 由数据卡重新计算月球表面的
        重力加速度。
  \item 在地球上，一次良好的蹬地可使你的质心升高
        \qty{0.50}{m}。对同样的蹬地，所达高度
        与 $g$ 成反比（该关系将在
        \cref{ch:g10:energy-conservation}中证明）。同样的起跳
        在月球上能使你升高多少？
  \item 水星：$M = \qty{3.30e23}{kg}$，$R = \qty{2.44e6}{m}$。
        求其表面的重力加速度，并与火星比较。
        水星比火星小得多 --- 两者的数值为何
        如此接近？
  \item 在距地表多高处，重力加速度
        降为地表值的一半？
  \item 距地心多远处，地球对每千克物体的引力
        恰好降为月球表面的数值
        \qty{1.62}{N/kg}？用地球半径表示。
\end{enumerate}

\medskip
\noindent\textbf{第四部分 --- 称量太阳。}
地球沿近似圆形的轨道运行，轨道半径为
$d = \qty{1.496e11}{m}$，一年即 $T = \qty{3.156e7}{s}$。取以下结论为已知
（其证明见\cref{ch:g12:satellites-kepler}）：沿半径为
$d$ 的圆周运动、速率为 $v$ 的物体，需要指向圆心的引力
为每千克 $v^2/d$ 牛顿。
\begin{enumerate}[resume]
  \item 计算地球的公转速率 $v$。
  \item 推算太阳施于地球的每千克引力。
  \item 该引力也等于 $G M_{\text{太阳}} / d^2$，由此推算
        太阳的质量。
  \item 相当于多少个地球？比较 $M_{\text{太阳}}$ 与
        $M_{\text{地球}}$。
  \item 尾声。卡文迪什的铅球、地球与太阳：一个常量
        $G$ 同时适用于三者。用一句话说明“万有引力”中的
        \emph{普遍}一词在
        本题中带来了什么好处。
\end{enumerate}
\end{problem}
